\clearpage\pagestyle{otMain}\setlength{\otBlocEspacement}{2.5mm}\fontsize{11}{15}\selectfont\corrigesChapitre{10}{Statistiques et diagrammes}
\begin{blocLivre}{cor-10-1}\begin{corrige}{1}{Un effectif total}\demarche{Additionner les effectifs avant de calculer la fréquence.}
Le total vaut $3+5+2=\res{10}$. La fréquence de 10 est $5/10=\res{0{,}5}$, soit 50\,\%.
\end{corrige}\end{blocLivre}
\begin{blocLivre}{cor-10-2}\begin{corrige}{2}{Une moyenne simple}\demarche{Additionner toutes les valeurs et compter les données.}
La somme vaut $6+9+10+15=40$. On divise par quatre valeurs : $40/4=\res{10}$.
\end{corrige}\end{blocLivre}
\begin{blocLivre}{cor-10-3}\begin{corrige}{3}{Une médiane impaire}\demarche{Ranger avant de chercher le rang central.}
La série rangée est $7,8,9,12,15$. L’effectif est 5 : la 3e valeur est centrale. La médiane vaut $\res9$.
\end{corrige}\end{blocLivre}
\begin{blocLivre}{cor-10-4}\begin{corrige}{4}{Une médiane paire}\demarche{Repérer les deux rangs centraux.}
Il y a six valeurs déjà rangées. Les valeurs centrales sont aux rangs 3 et 4 : 8 et 10. Leur moyenne vaut $\res9$. Trois données sont au plus 9 et trois sont au moins 9.
\end{corrige}\end{blocLivre}
\begin{blocLivre}{cor-10-5}\begin{corrige}{5}{Un secteur circulaire}\demarche{Prendre la fraction correspondante de 360 degrés.}
$30/100\times360=\res{108^\circ}$. L’angle est proportionnel à la part du total.
\end{corrige}\end{blocLivre}
\begin{blocLivre}{cor-10-6}\begin{corrige}{6}{Une moyenne avec effectifs}\demarche{Tenir compte des répétitions.}
L’effectif total vaut 10. La somme pondérée est $2\times4+5\times3+8\times3=47$. La moyenne est $47/10=\res{4{,}7}$. Il ne faut pas moyenner seulement 2, 5 et 8.
\end{corrige}\end{blocLivre}
\begin{blocLivre}{cor-10-7}\begin{corrige}{7}{Lire un diagramme en bâtons}\demarche{Cumuler les effectifs pour retrouver les rangs.}
Les effectifs sont 2 pour 0 livre, 4 pour 1, 3 pour 2, 2 pour 3 et 1 pour 4. Les rangs 6 et 7 valent respectivement 1 et 2. La médiane choisie est donc $(1+2)/2=\res{1{,}5}$ livre.
\label{fig-cor-10-7}\figureGeo[corrige]{stat-livres}
\end{corrige}\end{blocLivre}
\begin{blocLivre}{cor-10-8}\begin{corrige}{8}{Construire le partage}\demarche{Transformer chaque fréquence en angle.}
Les fréquences sont $1/2$, $3/10$, $1/5$. Les angles correspondants sont $\res{180^\circ,\ 108^\circ,\ 72^\circ}$. Leur somme vaut 360 degrés. On construit les secteurs à la suite et on ajoute une légende.
\label{fig-cor-10-8}\figureGeo[corrige]{stat-sports}
\end{corrige}\end{blocLivre}
\begin{blocLivre}{cor-10-9}\begin{corrige}{9}{Lire un pourcentage}\demarche{Passer de l’angle à la fréquence puis à l’effectif.}
La fréquence de A vaut $144/360=0{,}4=40\,\%$. Son effectif vaut donc $75\times0{,}4=\res{30}$.
\label{fig-cor-10-9}\figureGeo[corrige]{stat-secteur-a}
\end{corrige}\end{blocLivre}
\begin{blocLivre}{cor-10-10}\begin{corrige}{10}{Une formule de tableur}\demarche{Construire les produits avant de les additionner.}
En C2, saisir $\res{\texttt{=A2*B2}}$ puis recopier jusqu’à C4. La moyenne se calcule avec $\res{\texttt{=SOMME(C2:C4)/SOMME(B2:B4)}}$. Les produits valeur-effectif représentent la somme des données répétées.
\end{corrige}\end{blocLivre}
\begin{blocLivre}{cor-10-11}\begin{corrige}{11}{Ajouter une valeur}\demarche{Retrouver la somme à partir de la moyenne.}
La somme initiale vaut $5\times12=60$. La nouvelle somme vaut $60+18=78$, pour 6 élèves. La moyenne est $78/6=\res{13}$. Moyenner 12 et 18 donnerait un mauvais poids aux cinq premiers élèves.
\end{corrige}\end{blocLivre}
\begin{blocLivre}{cor-10-12}\begin{corrige}{12}{Une valeur manquante}\demarche{Utiliser la moyenne pour retrouver le total.}
La somme doit valoir $4\times12=48$. Donc $x=48-(8+11+14)=\res{15}$. La série rangée est $8,11,14,15$ et la médiane choisie vaut $(11+14)/2=\res{12{,}5}$.
\end{corrige}\end{blocLivre}
\begin{blocLivre}{cor-10-13}\begin{corrige}{13}{Une valeur extrême}\demarche{Calculer séparément les deux indicateurs.}
Pour A, moyenne et médiane valent 6. Pour B, la somme vaut 60 donc la moyenne vaut 12 ; la valeur centrale reste 6. La seule modification du dernier nombre change fortement la moyenne, sans changer la médiane.
\end{corrige}\end{blocLivre}
\begin{blocLivre}{cor-10-14}\begin{corrige}{14}{Comparer deux groupes}\demarche{Comparer des fréquences plutôt que des nombres bruts.}
A a une proportion $18/24=75\,\%$ et B $21/30=70\,\%$. A a donc la plus forte proportion, même si B compte davantage de sportifs. Les effectifs totaux sont différents.
\end{corrige}\end{blocLivre}
\begin{blocLivre}{cor-10-15}\begin{corrige}{15}{Interpréter une médiane}\demarche{Distinguer valeur centrale et fréquence d’une valeur.}
Non : une médiane parle de position dans la série, pas du nombre de valeurs égales à 15. Par exemple, dix durées de 10 min et dix de 20 min donnent une médiane $(10+20)/2=15$ min, sans aucune durée égale à 15.
\end{corrige}\end{blocLivre}
